% GENERATED FILE. Edit canonical lessons, then run npm run generate:books.
% canonical-source-hash: fnv1a64:aa0dadb37204349f
% canonical-lessons: ML-C356-sambhasanam, ML-C356-vivaranam, ML-C356-drsyam, ML-C356-kettariyuka, ML-C356-ketticcamaykkuka, ML-R356-review-349-356

\chapter{Telling It Again}
\label{ch:ml-telling-it-again}
\hlchaptermodality{356}

\begin{chapteropening}
\textbf{By the end of this chapter:} \emph{I can say a conversation, a description, footage, to learn by hearsay and to fabricate.}

Complete the last lesson of chapter 356: I can say a conversation, a description, footage, to learn by hearsay and to fabricate.
\end{chapteropening}

\section[\texorpdfstring{\ml{സംഭാഷണം} (sambhāṣaṇaṁ)}{സംഭാഷണം (sambhāṣaṇaṁ)}]{\ml{സംഭാഷണം} (sambhāṣaṇaṁ) — a conversation; dialogue}

\label{lesson:ML-C356-sambhasanam}



\begin{quote}
Before the new one: say the Malayalam for to gather, then the Malayalam for to boo.
\end{quote}

\subsection*{You'll want to know: \ml{സംഭാഷണം}}
\textbf{\ml{സംഭാഷണം}} — \emph{sambhāṣaṇaṁ} — \textquotedblleft{}a conversation; dialogue\textquotedblright{}.

\begin{quote}
\ml{ആ} \ml{സിനിമയിലെ} \ml{സംഭാഷണം} \ml{വളരെ} \ml{രസകരമായിരുന്നു}.
\end{quote}

\begin{quote}
\emph{ā sinimayile sambhāṣaṇaṁ vaḷare rasakaramāyirunnu.}
\end{quote}

\begin{quote}
The dialogue in that film was very interesting.
\end{quote}

The written word for a conversation, and the dialogue of a film or a play.

\begin{cousinweb}
\textbf{\ml{സംഭാഷണം}} (\emph{sambhāṣaṇaṁ}) is borrowed from Sanskrit \emph{sambhāṣaṇa}, \textquotedblleft{}conversation, discourse with someone\textquotedblright{} (Monier-Williams), built on \emph{bhāṣ}, \textquotedblleft{}to speak\textquotedblright{}, the root of \ml{ഭാഷ} (\emph{bhāṣa}), \textquotedblleft{}language\textquotedblright{}.
\end{cousinweb}

\begin{grammarlens}[title={What you've built}]
One more word for telling what happened.
\end{grammarlens}

\subsection*{Your turn}
\begin{itemize}
\raggedright
  \item \emph{Say it:} \emph{sambhāṣaṇaṁ}
  \item \emph{Say it:} \emph{sambhāṣaṇaṁ}, once more
  \item \emph{Say it:} \emph{ā sinimayile sambhāṣaṇaṁ vaḷare rasakaramāyirunnu.} — The dialogue in that film was very interesting.
  \item \emph{Recall:} say the Malayalam for to gather, then the Malayalam for to boo, then say \emph{sambhāṣaṇaṁ} again
\end{itemize}

\begin{tcolorbox}[breakable,colback=teal!4,colframe=teal!35!black,title={Before you move on}]
What does \ml{സംഭാഷണം} mean? (\textquotedblleft{}A conversation; dialogue\textquotedblright{}.) Say it once more.
\end{tcolorbox}
\section[\texorpdfstring{\ml{വിവരണം} (vivaraṇaṁ)}{വിവരണം (vivaraṇaṁ)}]{\ml{വിവരണം} (vivaraṇaṁ) — a description; an account}

\label{lesson:ML-C356-vivaranam}



\begin{quote}
Before the new one: say the Malayalam for to boo, then the Malayalam for a conversation.
\end{quote}

\subsection*{You'll want to know: \ml{വിവരണം}}
\textbf{\ml{വിവരണം}} — \emph{vivaraṇaṁ} — \textquotedblleft{}a description; an account\textquotedblright{}.

\begin{quote}
\ml{സാക്ഷി} \ml{സംഭവത്തിന്റെ} \ml{വിവരണം} \ml{പോലീസിനോട്} \ml{പറഞ്ഞു}.
\end{quote}

\begin{quote}
\emph{sākṣi sambhavattinṟe vivaraṇaṁ pōlīsinōṭŭ paṟaññu.}
\end{quote}

\begin{quote}
The witness gave the police an account of the incident.
\end{quote}

The noun of \ml{വിവരിക്കുക}, to describe: what you give when you describe.

\begin{cousinweb}
\textbf{\ml{വിവരണം}} (\emph{vivaraṇaṁ}) is borrowed from Sanskrit \emph{vivaraṇa}, \textquotedblleft{}uncovering, laying open; explanation\textquotedblright{} (Monier-Williams). It is the noun beside \ml{വിവരിക്കുക} (\emph{vivarikkuka}), \textquotedblleft{}to describe\textquotedblright{}.
\end{cousinweb}

\begin{grammarlens}[title={What you've built}]
One more word for telling what happened.
\end{grammarlens}

\subsection*{Your turn}
\begin{itemize}
\raggedright
  \item \emph{Say it:} \emph{vivaraṇaṁ}
  \item \emph{Say it:} \emph{vivaraṇaṁ}, once more
  \item \emph{Say it:} \emph{sākṣi sambhavattinṟe vivaraṇaṁ pōlīsinōṭŭ paṟaññu.} — The witness gave the police an account of the incident.
  \item \emph{Recall:} say the Malayalam for to boo, then the Malayalam for a conversation, then say \emph{vivaraṇaṁ} again
\end{itemize}

\begin{tcolorbox}[breakable,colback=teal!4,colframe=teal!35!black,title={Before you move on}]
What does \ml{വിവരണം} mean? (\textquotedblleft{}A description; an account\textquotedblright{}.) Say it once more.
\end{tcolorbox}
\section[\texorpdfstring{\ml{ദൃശ്യം} (dṛśyaṁ)}{ദൃശ്യം (dṛśyaṁ)}]{\ml{ദൃശ്യം} (dṛśyaṁ) — footage; a scene}

\label{lesson:ML-C356-drsyam}



\begin{quote}
Before the new one: say the Malayalam for a conversation, then the Malayalam for a description.
\end{quote}

\subsection*{You'll want to know: \ml{ദൃശ്യം}}
\textbf{\ml{ദൃശ്യം}} — \emph{dṛśyaṁ} — \textquotedblleft{}footage; a scene\textquotedblright{}.

\begin{quote}
\ml{അപകടത്തിന്റെ} \ml{ദൃശ്യം} \ml{ടിവിയിൽ} \ml{വന്നു}.
\end{quote}

\begin{quote}
\emph{apakaṭattinṟe dṛśyaṁ ṭiviyil vannu.}
\end{quote}

\begin{quote}
Footage of the accident came on TV.
\end{quote}

Plural: \ml{ദൃശ്യങ്ങൾ}, \emph{dṛśyaṅṅaḷ}.

Plural \ml{ദൃശ്യങ്ങൾ}: the pictures or footage of an event.

\begin{cousinweb}
\textbf{\ml{ദൃശ്യം}} (\emph{dṛśyaṁ}) is borrowed from Sanskrit \emph{dṛśya}, \textquotedblleft{}visible, worth seeing\textquotedblright{} (Monier-Williams), from \emph{dṛś}, \textquotedblleft{}to see\textquotedblright{}.
\end{cousinweb}

\begin{grammarlens}[title={What you've built}]
One more word for telling what happened.
\end{grammarlens}

\subsection*{Your turn}
\begin{itemize}
\raggedright
  \item \emph{Say it:} \emph{dṛśyaṁ}
  \item \emph{Say it:} \emph{dṛśyaṁ}, once more
  \item \emph{Say it:} \emph{apakaṭattinṟe dṛśyaṁ ṭiviyil vannu.} — Footage of the accident came on TV.
  \item \emph{Recall:} say the Malayalam for a conversation, then the Malayalam for a description, then say \emph{dṛśyaṁ} again
\end{itemize}

\begin{tcolorbox}[breakable,colback=teal!4,colframe=teal!35!black,title={Before you move on}]
What does \ml{ദൃശ്യം} mean? (\textquotedblleft{}Footage; a scene\textquotedblright{}.) Say it once more.
\end{tcolorbox}
\section[\texorpdfstring{\ml{കേട്ടറിയുക} (kēṭṭaṟiyuka)}{കേട്ടറിയുക (kēṭṭaṟiyuka)}]{\ml{കേട്ടറിയുക} (kēṭṭaṟiyuka) — to learn by hearsay; to hear of}

\label{lesson:ML-C356-kettariyuka}



\begin{quote}
Before the new one: say the Malayalam for a description, then the Malayalam for footage.
\end{quote}

\subsection*{You'll want to know: \ml{കേട്ടറിയുക}}
\textbf{\ml{കേട്ടറിയുക}} — \emph{kēṭṭaṟiyuka} — \textquotedblleft{}to learn by hearsay; to hear of\textquotedblright{}.

\begin{quote}
\ml{ആ} \ml{സംഭവം} \ml{ഞാൻ} \ml{കേട്ടറിഞ്ഞു}.
\end{quote}

\begin{quote}
\emph{ā sambhavaṁ ñān kēṭṭaṟiññu.}
\end{quote}

\begin{quote}
I heard about that incident from others.
\end{quote}

Past: \ml{കേട്ടറിഞ്ഞു}, \emph{kēṭṭaṟiññu}.

Joined from \ml{കേട്ട്}, having heard, and \ml{അറിയുക}, to know: to know something only because you heard it.

\begin{cousinweb}
\textbf{\ml{കേട്ടറിയുക}} (\emph{kēṭṭaṟiyuka}) is a native Dravidian compound of \ml{കേട്ട്} (\emph{kēṭṭŭ}), \textquotedblleft{}having heard\textquotedblright{}, from \ml{കേൾക്കുക} (DEDR 2017), and \ml{അറിയുക} (\emph{aṟiyuka}), \textquotedblleft{}to know\textquotedblright{} (DEDR 314). What you know this way, you did not see yourself.
\end{cousinweb}

\begin{grammarlens}[title={What you've built}]
One more word for telling what happened.
\end{grammarlens}

\subsection*{Your turn}
\begin{itemize}
\raggedright
  \item \emph{Say it:} \emph{kēṭṭaṟiyuka}
  \item \emph{Say it:} \emph{kēṭṭaṟiyuka}, once more
  \item \emph{Say it:} \emph{ā sambhavaṁ ñān kēṭṭaṟiññu.} — I heard about that incident from others.
  \item \emph{Recall:} say the Malayalam for a description, then the Malayalam for footage, then say \emph{kēṭṭaṟiyuka} again
\end{itemize}

\begin{tcolorbox}[breakable,colback=teal!4,colframe=teal!35!black,title={Before you move on}]
What does \ml{കേട്ടറിയുക} mean? (\textquotedblleft{}To learn by hearsay; to hear of\textquotedblright{}.) Say it once more.
\end{tcolorbox}
\section[\texorpdfstring{\ml{കെട്ടിച്ചമയ്ക്കുക} (keṭṭiccamaykkuka)}{കെട്ടിച്ചമയ്ക്കുക (keṭṭiccamaykkuka)}]{\ml{കെട്ടിച്ചമയ്ക്കുക} (keṭṭiccamaykkuka) — to fabricate; to make up}

\label{lesson:ML-C356-ketticcamaykkuka}



\begin{quote}
Before the new one: say the Malayalam for footage, then the Malayalam for to learn by hearsay.
\end{quote}

\subsection*{You'll want to know: \ml{കെട്ടിച്ചമയ്ക്കുക}}
\textbf{\ml{കെട്ടിച്ചമയ്ക്കുക}} — \emph{keṭṭiccamaykkuka} — \textquotedblleft{}to fabricate; to make up\textquotedblright{}.

\begin{quote}
\ml{അവൻ} \ml{ആ} \ml{കഥ} \ml{മുഴുവൻ} \ml{കെട്ടിച്ചമച്ചു}.
\end{quote}

\begin{quote}
\emph{avan ā katha muḻuvan keṭṭiccamaccu.}
\end{quote}

\begin{quote}
He made up the whole story.
\end{quote}

Past: \ml{കെട്ടിച്ചമച്ചു}, \emph{keṭṭiccamaccu}.

Said of false stories and false evidence: \ml{കെട്ടിച്ചമച്ച} \ml{കഥ}, a made-up story.

\begin{cousinweb}
\textbf{\ml{കെട്ടിച്ചമയ്ക്കുക}} (\emph{keṭṭiccamaykkuka}) is a native Dravidian compound of \ml{കെട്ടുക} (\emph{keṭṭuka}), \textquotedblleft{}to tie, to build\textquotedblright{} (DEDR 1147), and \ml{ചമയ്ക്കുക} (\emph{camaykkuka}), \textquotedblleft{}to make, to put together\textquotedblright{} (DEDR 2342): a made-up story is tied and built out of nothing.
\end{cousinweb}

\begin{grammarlens}[title={What you've built}]
One more word for telling what happened.
\end{grammarlens}

\subsection*{Your turn}
\begin{itemize}
\raggedright
  \item \emph{Say it:} \emph{keṭṭiccamaykkuka}
  \item \emph{Say it:} \emph{keṭṭiccamaykkuka}, once more
  \item \emph{Say it:} \emph{avan ā katha muḻuvan keṭṭiccamaccu.} — He made up the whole story.
  \item \emph{Recall:} say the Malayalam for footage, then the Malayalam for to learn by hearsay, then say \emph{keṭṭiccamaykkuka} again
\end{itemize}

\begin{tcolorbox}[breakable,colback=teal!4,colframe=teal!35!black,title={Before you move on}]
What does \ml{കെട്ടിച്ചമയ്ക്കുക} mean? (\textquotedblleft{}To fabricate; to make up\textquotedblright{}.) Say it once more.
\end{tcolorbox}
\section[(dialogue)]{Words for telling what happened}

\label{lesson:ML-R356-review-349-356}



\begin{quote}
Say the Malayalam for a witness and a war.
\end{quote}

\subsection*{Your turn}
Say these aloud:

\begin{itemize}
\raggedright
  \item a witness, a war, a correspondent, to happen, to announce
  \item a headline, an editor, an edition, a reader, to print
  \item a broadcast, a listener, a reporter, to depict, to broadcast
  \item a press conference, a statement, the real facts, to exaggerate, to say sorry
  \item a rumour, speculation, fake, to circulate, to spread around
  \item a speech, a medium, to give a speech, to applaud, to mention
  \item a journalist, an allegation, to reveal, to gather, to boo
  \item a conversation, a description, footage, to learn by hearsay, to fabricate
\end{itemize}

\begin{tcolorbox}[breakable,colback=teal!4,colframe=teal!35!black,title={Before you move on}]
A witness? (\textbf{\ml{സാക്ഷി}}, \emph{sākṣi}.) A war? (\textbf{\ml{യുദ്ധം}}, \emph{yuddhaṁ}.) A correspondent? (\textbf{\ml{ലേഖകൻ}}, \emph{lēkhakan}.) To happen? (\textbf{\ml{സംഭവിക്കുക}}, \emph{sambhavikkuka}.) To announce? (\textbf{\ml{പ്രഖ്യാപിക്കുക}}, \emph{prakhyāpikkuka}.) A headline? (\textbf{\ml{തലക്കെട്ട്}}, \emph{talakkeṭṭŭ}.) An editor? (\textbf{\ml{പത്രാധിപർ}}, \emph{patrādhipar}.) An edition? (\textbf{\ml{പതിപ്പ്}}, \emph{patippŭ}.) A reader? (\textbf{\ml{വായനക്കാരൻ}}, \emph{vāyanakkāran}.) To print? (\textbf{\ml{അച്ചടിക്കുക}}, \emph{accaṭikkuka}.) A broadcast? (\textbf{\ml{സംപ്രേഷണം}}, \emph{samprēṣaṇaṁ}.) A listener? (\textbf{\ml{ശ്രോതാവ്}}, \emph{śrōtāvŭ}.) A reporter? (\textbf{\ml{റിപ്പോർട്ടർ}}, \emph{ṟippōrṭṭar}.) To depict? (\textbf{\ml{ചിത്രീകരിക്കുക}}, \emph{citrīkarikkuka}.) To broadcast? (\textbf{\ml{സംപ്രേഷണം} \ml{ചെയ്യുക}}, \emph{samprēṣaṇaṁ ceyyuka}.) A press conference? (\textbf{\ml{പത്രസമ്മേളനം}}, \emph{patrasammēḷanaṁ}.) A statement? (\textbf{\ml{പ്രസ്താവന}}, \emph{prastāvana}.) The real facts? (\textbf{\ml{സത്യാവസ്ഥ}}, \emph{satyāvastha}.) To exaggerate? (\textbf{\ml{പെരുപ്പിക്കുക}}, \emph{peruppikkuka}.) To say sorry? (\textbf{\ml{മാപ്പ്} \ml{പറയുക}}, \emph{māppŭ paṟayuka}.) A rumour? (\textbf{\ml{കിംവദന്തി}}, \emph{kiṁvadanti}.) Speculation? (\textbf{\ml{അഭ്യൂഹം}}, \emph{abhyūhaṁ}.) Fake? (\textbf{\ml{വ്യാജ}}, \emph{vyāja}.) To circulate? (\textbf{\ml{പ്രചരിക്കുക}}, \emph{pracarikkuka}.) To spread around? (\textbf{\ml{പ്രചരിപ്പിക്കുക}}, \emph{pracarippikkuka}.) A speech? (\textbf{\ml{പ്രസംഗം}}, \emph{prasaṁgaṁ}.) A medium? (\textbf{\ml{മാധ്യമം}}, \emph{mādhyamaṁ}.) To give a speech? (\textbf{\ml{പ്രസംഗിക്കുക}}, \emph{prasaṁgikkuka}.) To applaud? (\textbf{\ml{കയ്യടിക്കുക}}, \emph{kayyaṭikkuka}.) To mention? (\textbf{\ml{പരാമർശിക്കുക}}, \emph{parāmarśikkuka}.) A journalist? (\textbf{\ml{പത്രപ്രവർത്തകൻ}}, \emph{patrapravarttakan}.) An allegation? (\textbf{\ml{ആരോപണം}}, \emph{ārōpaṇaṁ}.) To reveal? (\textbf{\ml{വെളിപ്പെടുത്തുക}}, \emph{veḷippeṭuttuka}.) To gather? (\textbf{\ml{ഒത്തുകൂടുക}}, \emph{ottukūṭuka}.) To boo? (\textbf{\ml{കൂവുക}}, \emph{kūvuka}.) A conversation? (\textbf{\ml{സംഭാഷണം}}, \emph{sambhāṣaṇaṁ}.) A description? (\textbf{\ml{വിവരണം}}, \emph{vivaraṇaṁ}.) Footage? (\textbf{\ml{ദൃശ്യം}}, \emph{dṛśyaṁ}.) To learn by hearsay? (\textbf{\ml{കേട്ടറിയുക}}, \emph{kēṭṭaṟiyuka}.) To fabricate? (\textbf{\ml{കെട്ടിച്ചമയ്ക്കുക}}, \emph{keṭṭiccamaykkuka}.)
\end{tcolorbox}
